\section{Infrared Spectrum and Propagating Degrees of Freedom}
  \label{sec:ir_spectrum}

  \subsection{Linearization Around Minkowski Spacetime}
    \label{subsec:linearization-around-minkowski-spacetime}

    We consider perturbations about flat spacetime,
    \[
      g_{\mu\nu}
      =
      \eta_{\mu\nu}
      +
      h_{\mu\nu},
    \]
    and work to quadratic order in $h_{\mu\nu}$.
    In the infrared regime
    $R\ell_\chi^2 \ll 1$,
    the quadratic operator derived from
    $\delta^2 S_\Pi^{(L)}$
    takes the local form
    \[
      \mathcal{O}_{\mu\nu}^{\ \ \rho\sigma}
      =
      c_{\mathrm{EH}}\Box
      +
      \beta \Box^2
      +
      \mathcal{O}(\ell_\chi^4 \Box^3),
    \]
    where $c_{\mathrm{EH}}>0$ and $\beta<0$ in the natural spectral scheme.

    We impose harmonic gauge,
    \[
      \partial^\mu \bar h_{\mu\nu}
      =
      0,
      \qquad
      \bar h_{\mu\nu}
      =
      h_{\mu\nu}
      -
      \frac12 \eta_{\mu\nu} h,
    \]
    which eliminates gauge redundancies at linear order.

  \subsection{Scalar--Vector--Tensor Decomposition}
    \label{subsec:scalar--vector--tensor-decomposition}

    To identify the propagating modes, we perform a
    scalar--vector--tensor (SVT) decomposition with respect to spatial
    rotations.
    Writing spatial indices as $i,j=1,2,3$, the perturbation splits as
    \[
      h_{00},
      \qquad
      h_{0i} = \partial_i B + B_i,
      \qquad
      h_{ij}
      =
      2\psi\,\delta_{ij}
      +
      2\partial_i\partial_j E
      +
      \partial_i E_j
      +
      \partial_j E_i
      +
      h^{\mathrm{TT}}_{ij},
    \]
    with
    \[
      \partial_i B_i = 0,
      \qquad
      \partial_i E_i = 0,
      \qquad
      \partial_i h^{\mathrm{TT}}_{ij} = 0,
      \qquad
      h^{\mathrm{TT}}_{ii} = 0.
    \]

    In harmonic gauge, the scalar and vector components enter the quadratic
    action without independent kinetic terms in the regime
    $R\ell_\chi^2 \ll 1$.
    They are therefore non-dynamical and can be eliminated algebraically.

    The only sector acquiring a propagating kinetic operator is the
    transverse-traceless tensor sector $h^{\mathrm{TT}}_{ij}$.

  \subsection{Transverse-Traceless Sector}
    \label{subsec:transverse-traceless-sector}

    Restricting to $h^{\mathrm{TT}}_{\mu\nu}$,
    the quadratic operator reduces to
    \[
      \left(
        c_{\mathrm{EH}}\Box
        +
        \beta \Box^2
      \right)
      h^{\mathrm{TT}}_{\mu\nu}
      =
      0.
    \]

    Factoring gives
    \[
      \Box
      \left(
        1
        +
        \frac{\beta}{c_{\mathrm{EH}}}
        \Box
      \right)
      h^{\mathrm{TT}}_{\mu\nu}
      =
      0.
    \]

    The first factor corresponds to the standard massless graviton pole
    $k^2=0$.
    The second factor introduces an additional pole at
    \[
      k^2
      =
      -\frac{c_{\mathrm{EH}}}{\beta}
      \sim
      \ell_\chi^{-2}.
    \]

    Since $\ell_\chi^{-1}$ sets the pre-geometric ultraviolet scale,
    this excitation lies outside the infrared window
    \[
      \omega \ll \ell_\chi^{-1},
    \]
    relevant to gravitational-wave observations.

  \subsection{Infrared Propagation Theorem}
    \label{subsec:infrared-propagation-theorem}

    \paragraph{Proposition.}
      In the regime $R\ell_\chi^2 \ll 1$ and $\omega \ll \ell_\chi^{-1}$,
      the Lorentzian spectral action propagates exactly two transverse-traceless tensor modes of helicity $\pm 2$.
      No scalar or vector propagating degree of freedom is dynamically accessible in this domain.

    \paragraph{Proof.}
      The SVT decomposition shows that scalar and vector components lack
      independent kinetic operators in harmonic gauge.
      The tensor sector obeys a fourth-order equation that factorizes into two second-order operators.
      The additional pole is located at $k^2 \sim \ell_\chi^{-2}$, hence outside the infrared regime.
      Only the massless pole contributes to low-frequency propagation.
      \hfill $\square$

\subsection{Polarization Content}
  \label{subsec:polarization-content}

  The surviving infrared excitations are precisely the two helicity $\pm 2$ tensor modes of general relativity.
  No scalar (breathing or longitudinal) or vector polarizations appear in the accessible frequency range.
  The Lorentzian completion of the spectral framework is therefore
  consistent with current gravitational-wave polarization constraints.
